%%%%%%%%%%%%%%%%%%%%%%%%%%%%%%%%%%%%%%%%%%%%%%%%%%

%% CV-LuaLaTeX in English
%% Carla Cisternas
%% https://github.com/carlacisternasg/CV-LuaLaTeX

%% Based on the Awesome CV LaTeX Template
%% https://github.com/darwiin/awesome-neue-latex-cv
%% and the Bastián Gozález-Bustamante's CV-LaTeX Template
%% https://github.com/bgonzalezbustamante/CV-LaTeX

%% Creative Commons Attribution 4.0 International License
%% https://github.com/carlacisternasg/CV-LuaLaTeX/blob/master/LICENSE.md

%%%%%%%%%%%%%%%%%%%%%%%%%%%%%%%%%%%%%%%%%%%%%%%%%%

\par{I am a PhD Researcher in the Department of Latin American Studies at the Institute for History at the Faculty of Humanities at Leiden University, Netherlands. I hold an MA (1st) in Political Science from the Institute of Public Affairs at the Universidad de Chile. Moreover, I earned a BA (1st) and a BPS (1st) in Public Administration from the Universidad de Santiago de Chile (USACH).

My doctoral project is supervised by Prof. dr. Patricio Silva and Dr. María Gabriela Palacio and focuses on the relationship between technocracy, democracy, and sociopolitical conflicts. In particular, I study the role of external experts in the design of public policies in Chile between 1990 and 2022. My thesis follows an article-based format, the first of which, titled {\bfseries “Strategic Use of Ad Hoc Commissions for Blame Avoidance: Evidence from Chile,”} was published in {\itshape Public Administration Review}, one of the most prestigious journals in the field of public administration. The Chilean National Agency for Research and Development and the Slicher van Bath de Jong Foundation, through the Centre for Latin American Research and Documentation (CEDLA) of the University of Amsterdam, fund my PhD project.}\\